% mpe-beamer-theme.tex — thème Beamer reprenant mpe_pres_revealjs.css
% À inclure via  include-in-header: mpe-beamer-theme.tex

%% ---------- Couleurs (classes du CSS) ----------
\definecolor{jaune}{HTML}{F1D763}
\definecolor{jauneb}{HTML}{FDC455}
\definecolor{orange}{HTML}{F7A913}
\definecolor{rouge}{HTML}{C94326}
\definecolor{vert}{HTML}{1F908E}
\definecolor{bleu}{HTML}{33658A}
\definecolor{bleuf}{HTML}{0C324E}
\definecolor{clair}{HTML}{99B6B2}
\definecolor{alea}{HTML}{398D55}
\definecolor{fixe}{HTML}{33658A}
\definecolor{lien}{HTML}{B40F20}
\definecolor{care}{HTML}{B40F20}
\definecolor{highlight}{HTML}{B40F20}
\definecolor{calloutbg}{HTML}{F0F2F0}
\definecolor{tabgris}{HTML}{E4E8E4}

%% ---------- Polices : Latin Modern (fournie avec toute distribution TeX) ----------
% Texte en Latin Modern Sans, maths en Latin Modern classique.
% Polices chargées par nom de fichier : fonctionne aussi avec TinyTeX /
% dans un conteneur Docker, sans dépendre de fontconfig.
% Alternative sobre type Helvetica : décommenter le bloc TeX Gyre Heros.
\usepackage{iftex}
\ifPDFTeX\else
  \usepackage{fontspec}
  \setsansfont{lmsans10-regular.otf}[
    BoldFont       = lmsans10-bold.otf,
    ItalicFont     = lmsans10-oblique.otf,
    BoldItalicFont = lmsans10-boldoblique.otf]
  % \setsansfont{texgyreheros-regular.otf}[Scale=0.92,
  %   BoldFont = texgyreheros-bold.otf, ItalicFont = texgyreheros-italic.otf,
  %   BoldItalicFont = texgyreheros-bolditalic.otf]
  \setmonofont{lmmono10-regular.otf}[
    BoldFont = lmmonolt10-bold.otf, ItalicFont = lmmono10-italic.otf]
\fi
\usefonttheme[onlymath]{serif}

%% ---------- Couleurs Beamer ----------
\setbeamercolor{normal text}{fg=black!85}
\setbeamercolor{structure}{fg=bleu}
\setbeamercolor{frametitle}{fg=rouge,bg=}              % h2 du CSS
\setbeamercolor{title}{fg=bleu}
\setbeamercolor{subtitle}{fg=vert}
\setbeamercolor{author}{fg=bleuf}
\setbeamercolor{institute}{fg=vert}
\setbeamercolor{date}{fg=gray}
\setbeamercolor{section title}{fg=bleu}                % h1 du CSS
\setbeamercolor{block title}{fg=bleuf,bg=}             % h3 du CSS
\setbeamercolor{block body}{bg=}
\setbeamercolor{item}{fg=rouge}
\setbeamercolor{alerted text}{fg=lien}
\setbeamercolor{caption name}{fg=bleu}
\setbeamercolor{page number in head/foot}{fg=gray}

\setbeamerfont{frametitle}{size=\large}
\setbeamerfont{block title}{size=\normalsize,series=\bfseries}
\setbeamerfont{caption}{size=\scriptsize,shape=\itshape}

\setbeamertemplate{navigation symbols}{}
\setbeamertemplate{itemize items}[circle]
\setbeamertemplate{caption}[default]
\setbeamertemplate{footline}{%
  \hfill\usebeamercolor[fg]{page number in head/foot}%
  \scriptsize\insertframenumber\,/\,\inserttotalframenumber\hspace*{1.2em}%
  \vspace*{1.2ex}}

% Page de section (équivalent des slides "# Titre" en revealjs)
\setbeamertemplate{section page}{%
  \vspace*{\fill}
  \begin{center}
    {\usebeamercolor[fg]{section title}\usebeamerfont{title}\insertsectionhead\par}
    \vspace{1ex}{\color{orange}\rule{0.4\linewidth}{1.5pt}}
  \end{center}
  \vspace*{\fill}}

\AtBeginDocument{\hypersetup{colorlinks=true,urlcolor=lien,linkcolor=.}}

%% ---------- Callouts (callout-note & co) ----------
\usepackage{tcolorbox}
\tcbuselibrary{skins}
\tcbset{mpecalloutstyle/.style={enhanced, frame hidden, colback=calloutbg,
    arc=0pt, outer arc=0pt, left=6pt, right=6pt, top=4pt, bottom=4pt,
    before skip=4pt, after skip=6pt}}
\newenvironment{mpecallout}[2][rouge]{%
  \tcolorbox[mpecalloutstyle, borderline west={4pt}{0pt}{#1}]%
  \if\relax\detokenize{#2}\relax\else
    {\color{#1}\usebeamerfont{block title}#2\par}\smallskip
  \fi}{\endtcolorbox}

%% ---------- Barre d'onglets (panel-tabset) ----------
\newcommand{\qtab}[1]{{\setlength{\fboxsep}{2.5pt}%
  \colorbox{tabgris}{\strut\scriptsize\color{bleuf}#1}}}
\newcommand{\qtabactive}[1]{{\setlength{\fboxsep}{2.5pt}%
  \colorbox{rouge}{\strut\scriptsize\color{white}\bfseries #1}}}
\newcommand{\qtabsep}{\hspace{2pt}\allowbreak}
\newenvironment{qtabbar}{\par\vspace*{-0.8ex}\noindent\raggedright}%
  {\par\vspace{-0.9ex}\noindent{\color{rouge}\rule{\linewidth}{0.6pt}}\par\medskip}

%% ---------- Macros mathématiques (bloc ::: hidden de la version revealjs) ----------
\providecommand\R{{\mathbb{R}}}
\providecommand\Xbf{{\boldsymbol{X}}}
\providecommand\norm[1]{\lVert#1\rVert}
\providecommand\xcol[1]{\boldsymbol{x}^{#1}}
\providecommand\xrow[1]{\boldsymbol{x}_{#1}}
\providecommand\xbf{\boldsymbol{x}}
\providecommand\ybf{\boldsymbol{y}}
% Macros du cours de statistique spatiale
\providecommand\E{{\mathbb{E}}}
\renewcommand\P{{\mathbb{P}}}   % \P existe déjà en LaTeX (pied-de-mouche)
\providecommand\Cov{{\mathbb{C}\text{ov}}}
\providecommand\Var{{\mathbb{V}\text{ar}}}
\providecommand\Zchap{\widehat{Z}}
\providecommand\un{\mathbf{1}}
% \class{couleur}{...} (extension HTML de MathJax) -> texte coloré en LaTeX
\providecommand\class[2]{\textcolor{#1}{#2}}
% Macro de bordure de flextable : définie ici car sa définition (avec #1)
% dans une frame beamer fait planter LaTeX ; le filtre la retire des tableaux.
\providecommand{\ascline}[3]{\noalign{\global\arrayrulewidth #1}\arrayrulecolor[HTML]{#2}\cline{#3}}

%% ---------- Taille maximale des figures ----------
% Les figures sont réduites pour tenir dans \mpefigmaxh de la hauteur de texte
% (et dans la largeur de la colonne). Modifier la fraction si besoin, ou
% localement dans une slide :  `\renewcommand{\mpefigmaxh}{0.4}`{=latex}
\newcommand{\mpefigmaxh}{0.6}
\makeatletter
\AtBeginDocument{%
  \@ifundefined{pandocbounded}{}{%
    \renewcommand*\pandocbounded[1]{%
      \sbox\pandoc@box{#1}%
      \Gscale@div\@tempa{\mpefigmaxh\textheight}{\dimexpr\ht\pandoc@box+\dp\pandoc@box\relax}%
      \Gscale@div\@tempb{\linewidth}{\wd\pandoc@box}%
      \ifdim\@tempb\p@<\@tempa\p@\let\@tempa\@tempb\fi
      \ifdim\@tempa\p@<\p@\scalebox{\@tempa}{\usebox\pandoc@box}%
      \else\usebox{\pandoc@box}\fi}}%
  \@ifundefined{maxheight}{}{%
    \def\maxheight{\ifdim\Gin@nat@height>\mpefigmaxh\textheight
      \mpefigmaxh\textheight\else\Gin@nat@height\fi}}%
}
\makeatother
